\chapter{XDP --- \textit{eXpress Data Path}}

O \textit{eXpress Data Path} (XDP) é o ponto de intervenção mais precoce
disponível para programas eBPF na pilha de rede do kernel Linux \cite{linuxkernel2024}.
Ao executar diretamente no driver da interface de rede --- antes da alocação do
\texttt{sk\_buff} e antes de qualquer processamento pela pilha TCP/IP --- o XDP
permite processar pacotes com latência mínima e overhead desprezível em comparação
com abordagens em espaço de usuário \cite{cassagnes2020}.
Em ambientes de funções virtualizadas de rede, onde o desempenho de encaminhamento
e a capacidade de monitoramento sem intrusão são requisitos fundamentais, o XDP
representa uma das contribuições mais relevantes do eBPF \cite{venancio2021}.

% ----------------------------------------------------------------
\section{Modos de execução}
% ----------------------------------------------------------------

O XDP suporta três modos de execução, que diferem em como o programa eBPF é
integrado ao driver da interface de rede e, consequentemente, no desempenho
e na compatibilidade obtidos \cite{linuxkernel2024}.

\textbf{XDP Nativo (\textit{Native XDP}).}
O programa eBPF é executado diretamente no contexto do driver da interface de rede,
antes da alocação do \texttt{sk\_buff} \cite{linuxkernel2024}.
Este modo oferece o melhor desempenho possível, pois elimina toda a sobrecarga
de inicialização da pilha de rede para pacotes que serão descartados ou
redirecionados.
Requer suporte explícito no driver da interface de rede; a maioria dos drivers
de alto desempenho --- como \texttt{mlx5}, \texttt{i40e}, \texttt{ixgbe} e
\texttt{virtio\_net} --- já oferece esse suporte \cite{gregg2019}.

\textbf{XDP Genérico (\textit{Generic XDP}).}
Disponível em qualquer interface de rede, independentemente do suporte do driver
\cite{linuxkernel2024}.
Neste modo, o programa eBPF é executado após a alocação do \texttt{sk\_buff},
dentro do núcleo da pilha de rede, o que significa que o overhead de alocação
de memória já foi incorrido antes da execução do programa.
O desempenho é inferior ao XDP nativo, mas este modo é útil para desenvolvimento,
testes e implantação em interfaces sem suporte nativo, como as de máquinas
virtuais mais simples \cite{gregg2020}.

\textbf{XDP Offload (\textit{Offloaded XDP}).}
O programa eBPF é compilado e carregado diretamente no firmware da placa de rede,
que o executa sem qualquer envolvimento da CPU do host \cite{linuxkernel2024}.
Este modo oferece o máximo de desempenho e a menor latência possível, pois o
processamento ocorre inteiramente no hardware de rede.
Atualmente, o suporte é limitado a um pequeno conjunto de placas de rede de alta
performance com capacidade de processamento programável, como as da família
Netronome Agilio \cite{gregg2019}.

A verificação do modo suportado pela interface pode ser feita com:

\begin{lstlisting}[language=bash,
  caption={Verificação do modo XDP suportado pela interface de rede.}]
ip link show eth0
# Saida esperada para XDP nativo ativo:
# xdp/id:42  <BROADCAST,MULTICAST,UP,LOWER_UP>
\end{lstlisting}

Para fins didáticos, o minicurso utiliza o modo genérico em máquinas virtuais,
garantindo compatibilidade independentemente do driver de rede disponível
\cite{plinta2026repo}.

% ----------------------------------------------------------------
\section{XDP e processamento de pacotes}
% ----------------------------------------------------------------

Um programa XDP recebe como argumento um ponteiro para a estrutura
\texttt{struct xdp\_md}, que expõe os limites do buffer do pacote em memória
\cite{linuxkernel2024}:

\begin{lstlisting}[language=C,
  caption={Estrutura de contexto de um programa XDP.}]
struct xdp_md {
    __u32 data;          /* inicio dos dados do pacote  */
    __u32 data_end;      /* fim dos dados do pacote     */
    __u32 data_meta;     /* area de metadados           */
    __u32 ingress_ifindex;
    __u32 rx_queue_index;
};
\end{lstlisting}

O acesso aos bytes do pacote é feito por meio de ponteiros calculados a partir de
\texttt{data} e \texttt{data\_end}.
Toda operação de leitura ou escrita deve ser precedida de verificação explícita de
limites, que o verificador do kernel valida estaticamente \cite{linuxkernel2024}:

\begin{lstlisting}[language=C,
  caption={Acesso seguro ao cabeçalho Ethernet em um programa XDP.}]
void *data     = (void *)(long)ctx->data;
void *data_end = (void *)(long)ctx->data_end;

struct ethhdr *eth = data;
if ((void *)(eth + 1) > data_end)
    return XDP_PASS;   /* pacote muito curto; entregar a pilha */
\end{lstlisting}

Ao final do processamento, o programa retorna um código de ação que determina o
destino do pacote \cite{linuxkernel2024}:

\begin{itemize}
  \item \textbf{\texttt{XDP\_PASS}} --- entrega o pacote à pilha de rede normal
        do kernel para processamento convencional;
  \item \textbf{\texttt{XDP\_DROP}} --- descarta o pacote imediatamente, sem
        alocar \texttt{sk\_buff} nem percorrer a pilha TCP/IP; é a ação de menor
        custo computacional e a preferida para mitigação de DDoS \cite{gregg2019};
  \item \textbf{\texttt{XDP\_TX}} --- retransmite o pacote pela mesma interface
        de entrada, possivelmente após modificação de cabeçalhos;
  \item \textbf{\texttt{XDP\_REDIRECT}} --- encaminha o pacote para outra
        interface de rede, para uma CPU diferente ou diretamente para um socket
        em espaço de usuário via AF\_XDP \cite{linuxkernel2024};
  \item \textbf{\texttt{XDP\_ABORTED}} --- indica erro de programa; o pacote é
        descartado e um evento de rastreamento é gerado para depuração.
\end{itemize}

Para monitoramento de funções virtualizadas de rede, a ação \texttt{XDP\_PASS}
combinada com coleta de métricas em mapas BPF é o padrão mais utilizado:
o programa observa cada pacote --- contando bytes, registrando protocolos e
medindo intervalos entre chegadas --- sem interferir no fluxo normal de dados
\cite{venancio2021}.

O exemplo a seguir demonstra um programa XDP de monitoramento que acumula bytes
recebidos por protocolo IP em um mapa \texttt{PERCPU\_ARRAY} para exportação via
Prometheus:

\begin{lstlisting}[language=C,
  caption={Programa XDP de monitoramento de bytes por protocolo L4.}]
SEC("xdp")
int monitor_bytes(struct xdp_md *ctx)
{
    void *data     = (void *)(long)ctx->data;
    void *data_end = (void *)(long)ctx->data_end;

    struct ethhdr *eth = data;
    if ((void *)(eth + 1) > data_end)
        return XDP_PASS;

    if (eth->h_proto != bpf_htons(ETH_P_IP))
        return XDP_PASS;

    struct iphdr *ip = (void *)(eth + 1);
    if ((void *)(ip + 1) > data_end)
        return XDP_PASS;

    __u32 proto  = ip->protocol;
    __u32 nbytes = (void *)data_end - data;

    __u64 *cnt = bpf_map_lookup_elem(&bytes_by_proto, &proto);
    if (cnt)
        *cnt += nbytes;

    return XDP_PASS;
}
\end{lstlisting}

% ----------------------------------------------------------------
\section{Redirecionamento de pacotes}
% ----------------------------------------------------------------

O redirecionamento de pacotes com \texttt{XDP\_REDIRECT} é um dos casos de uso
mais poderosos do XDP para funções virtualizadas de rede, pois permite implementar
encaminhamento de alta velocidade inteiramente em eBPF, sem passar pela pilha
de roteamento do kernel \cite{linuxkernel2024}.

\textbf{Redirecionamento entre interfaces.}
A \textit{helper function} \texttt{bpf\_redirect} encaminha o pacote para uma
interface de saída identificada por seu índice (\texttt{ifindex}):

\begin{lstlisting}[language=C,
  caption={Redirecionamento de pacote para outra interface com XDP.}]
SEC("xdp")
int encaminhar(struct xdp_md *ctx)
{
    /* encaminha todos os pacotes para a interface de indice 3 */
    return bpf_redirect(3, 0);
}
\end{lstlisting}

Esse padrão é utilizado para implementar pontes (\textit{bridges}) e roteadores
de alto desempenho em software, substituindo soluções baseadas em \textit{iptables}
com ganhos expressivos de throughput \cite{gregg2019}.

\textbf{Redirecionamento para socket AF\_XDP.}
O AF\_XDP é uma família de sockets que permite que aplicações em espaço de usuário
recebam pacotes diretamente do driver de rede, sem passar pela pilha TCP/IP,
por meio de filas de memória compartilhada (\textit{UMEM}) entre kernel e espaço
de usuário \cite{linuxkernel2024}.
A \textit{helper function} \texttt{bpf\_redirect\_map} com um mapa do tipo
\texttt{BPF\_MAP\_TYPE\_XSKMAP} é usada para selecionar quais pacotes são
entregues a sockets AF\_XDP, enquanto os demais seguem o caminho normal com
\texttt{XDP\_PASS}:

\begin{lstlisting}[language=C,
  caption={Redirecionamento seletivo de pacotes UDP para socket AF\_XDP.}]
SEC("xdp")
int xsk_redirect(struct xdp_md *ctx)
{
    /* ... parse dos cabecalhos ... */
    if (ip->protocol == IPPROTO_UDP) {
        __u32 idx = ctx->rx_queue_index;
        return bpf_redirect_map(&xsk_map, idx, XDP_PASS);
    }
    return XDP_PASS;
}
\end{lstlisting}

Em ambientes NFV, o AF\_XDP é utilizado por funções de rede que precisam processar
pacotes em espaço de usuário com latência próxima ao kernel --- como firewalls e
balanceadores de carga implementados em \textit{Data Plane Development Kit} (DPDK)
--- mas com a vantagem de interoperabilidade com a pilha de rede padrão do Linux
para o tráfego não selecionado \cite{venancio2021}.

% ----------------------------------------------------------------
\section{Vantagens e limitações}
% ----------------------------------------------------------------

\textbf{Vantagens.}

\textbf{Desempenho.}
O XDP nativo é a abordagem de maior desempenho para processamento de pacotes em
software no Linux, superando soluções baseadas em \textit{iptables} e em módulos
de kernel convencionais em cenários de alta taxa de pacotes \cite{cassagnes2020}.
A ausência de alocação de \texttt{sk\_buff} para pacotes descartados ou
redirecionados elimina o principal gargalo de memória da pilha de rede em cargas
de trabalho de filtragem \cite{gregg2019}.

\textbf{Programabilidade dinâmica.}
Programas XDP podem ser carregados e substituídos em tempo de execução sem
interromper o tráfego de rede, ao contrário de filtros estáticos de
\textit{iptables} que exigem reinicialização do serviço de firewall
\cite{linuxkernel2024}.
Mapas BPF compartilhados entre o programa e o espaço de usuário permitem atualizar
configurações --- listas de endereços bloqueados, tabelas de roteamento,
limiares de monitoramento --- sem recarregar o programa \cite{gregg2020}.

\textbf{Visibilidade sem intrusão.}
A ação \texttt{XDP\_PASS} permite que programas XDP observem todo o tráfego de
uma interface e coletem métricas detalhadas sem alterar o comportamento da rede
--- propriedade fundamental para monitoramento de funções virtualizadas de rede
em produção \cite{venancio2021}.

\textbf{Limitações.}

\textbf{Acesso restrito ao contexto do pacote.}
No modo nativo, o programa XDP tem acesso apenas aos bytes do pacote via
\texttt{xdp\_md} --- sem acesso a metadados de cgroup, namespace ou PID do
processo originador, que só estão disponíveis em \textit{hooks} posteriores na
pilha, como o TC \cite{linuxkernel2024}.
Para monitoramento que exige associação de fluxos de rede a funções virtualizadas
específicas, o TC ou \textit{tracepoints} de socket oferecem visibilidade mais rica.

\textbf{Suporte limitado de drivers.}
O XDP nativo requer suporte explícito no driver da interface de rede
\cite{linuxkernel2024}.
Em ambientes de máquina virtual com drivers paravirtualizados menos comuns, pode
ser necessário recorrer ao XDP genérico, com desempenho inferior
\cite{gregg2020}.

\textbf{Ausência de estado de conexão.}
O XDP opera pacote a pacote e não mantém estado de conexão nativo
\cite{linuxkernel2024}.
Casos de uso que exigem rastreamento de fluxos TCP --- como detecção de
retransmissões ou medição de RTT --- precisam manter esse estado explicitamente
em mapas BPF, com o cuidado adicional de gerenciar a expiração de entradas para
evitar esgotamento de memória \cite{gregg2019}.

\textbf{Restrições do verificador.}
Como todo programa eBPF, programas XDP estão sujeitos às restrições do verificador
--- limite de complexidade de laços, verificação obrigatória de limites de acesso
a memória e ausência de recursão \cite{linuxkernel2024}.
Lógica de processamento muito complexa pode não ser verificável e precisar ser
distribuída entre múltiplos programas encadeados via \textit{tail calls}
\cite{gregg2019}.